\NSPage{049}%
equation across \NSi{NS-F-P049-001}{\xi=0}. The original equations preserve even parity of the first four coordinates
and odd parity of the last two; uniqueness gives precisely these parities. Taylor's formula for
an even smooth function of \NSi{NS-F-P049-002}{\xi} shows, to every finite order, that it is smooth as a function of
\NSi{NS-F-P049-003}{X=\xi^2} on the half-interval. This also verifies the regularity required to form the next source
\NSi{NS-F-P049-004}{\Omega_{n}/X}.
\end{proof}

\subsection{Canceling the total radial residual integrals}\label{sec:5.2}
The inner profiles \NSi{NS-F-P049-005}{(\phi_n,U_n,V_n,\Pi_n)} from
Lemma~\ref{lem:5.1} are not yet coefficients that can be used at the next order. We must extend
\NSi{NS-F-P049-006}{E_n,U_n} to all \NSi{NS-F-P049-007}{X\geq 0}, recompute \NSi{NS-F-P049-008}{V_n,\Pi_n} from
Equations~\eqref{eq:5.2} and~\eqref{eq:5.5}, and control the residual created by that extension.
A radial cutoff of \NSi{NS-F-P049-009}{E_n,U_n} alone does not give compact support for the fields obtained by radial
integration. In particular, the pressure may acquire a nonzero exterior constant and the stress
\NSi{NS-F-P049-010}{\mathcal T_n} defined below in \eqref{eq:5.9} may acquire a nonzero exterior tail. We impose
vanishing of the five moments defined below in Equations~\eqref{eq:5.10} and~\eqref{eq:5.11} to remove these
terms, as proved in Lemma~\ref{lem:5.2}.

We specify how to collect the residual coefficients before imposing the moments, using
\NSi{NS-F-P049-011}{R=\sqrt{2X}}. To extract order \NSi{NS-F-P049-012}{n}, we substitute \eqref{eq:5.1} into the differential
polynomial \NSi{NS-F-P049-013}{\mathscr R}, perform the physical differentiations, remove the common tangential power
\NSi{NS-F-P049-014}{q^{-A-1}}, and collect the coefficient of \NSi{NS-F-P049-015}{q^{2nh}}. This is a finite algebraic
operation: multiplication uses the sums \NSi{NS-F-P049-016}{i+j=n}, and axial viscosity uses order
\NSi{NS-F-P049-017}{n-1}. Explicitly,
\NSFormula{NS-F-P049-018}{display}{none}
\[
\begin{aligned}
\mathfrak r_{\theta,n}
  &=\frac{R}{\mathsf C}\left[
      T_{b,n}\phi_n
      +\sum_{i+j=n}\left\{V_i\left(\phi_j'+\frac{\phi_j}{X}\right)+U_iZ_{b,j}\phi_j\right\}
      -2\left(X\phi_n''+2\phi_n'\right)
      -Z_{b,n-1}^{[2]}\phi_{n-1}
    \right],\\
\mathfrak r_{z,n}
  &=T_{c,n}U_n
      +\sum_{i+j=n}\left\{V_iU_j'+U_iZ_{c,j}U_j\right\}
      +Z_{-2A,n}\Pi_n
      -2\left(XU_n''+U_n'\right)
      -Z_{c,n-1}^{[2]}U_{n-1}.
\end{aligned}
\]
These expressions vanish where the inner equations are solved. We define the two required stress
components by weighted radial integration of the negative residuals:
\NSFormula{NS-F-P049-019}{display}{5.9}
\begin{equation}\label{eq:5.9}
\mathcal T_{n,\theta}(R)=-R^{-2}\int_0^R\varrho^2\mathfrak r_{\theta,n}(\varrho)\,d\varrho,
\qquad
\mathcal T_{n,z}(R)=-R^{-1}\int_0^R\varrho\mathfrak r_{z,n}(\varrho)\,d\varrho.
\tag{5.9}
\end{equation}
All profiles under these integrals have the same parameter \NSi{NS-F-P049-020}{\eta}. The two physical stress
components are \NSi{NS-F-P049-021}{q^{-A-1/2+\lambda_n}\mathcal T_n}. For the forward integrals to vanish beyond the support
of their sources, the total integrals of \NSi{NS-F-P049-022}{R^2\mathfrak r_{\theta,n}} and
\NSi{NS-F-P049-023}{R\mathfrak r_{z,n}} must vanish. The five conditions below imply those two identities through
the conservative momentum equations; they also remove the radial-velocity and pressure terms left by
integration. We also write \NSi{NS-F-P049-024}{F_n=\int_0^XU_n(x,\eta)\,dx=X\mathcal A_X(U_n)} for the Stokes
streamfunction coefficient; this indexed notation is distinct from the leading azimuthal profile
\NSi{NS-F-P049-025}{F=E/\sqrt{2X}}.

For a tentative extension at order \NSi{NS-F-P049-026}{n}, we denote the five total moments by
\NSFormula{NS-F-P049-027}{display}{5.10}
\begin{equation}\label{eq:5.10}
m_{n,1}=\int_0^\infty RU_n\,dR,
\qquad
m_{n,2}=\int_0^\infty R^2E_n\,dR,
\qquad
m_{n,3}=\int_0^\infty \partial_R\Pi_n\,dR.
\tag{5.10}
\end{equation}
\NSFormula{NS-F-P049-028}{display}{5.11}
\begin{equation}\label{eq:5.11}
\begin{aligned}
m_{n,4}&=\int_0^\infty R^2\sum_{i+j=n}U_iE_j\,dR,\\
m_{n,5}&=\int_0^\infty\left(R\sum_{i+j=n}U_iU_j-\frac12R^2\partial_R\Pi_n\right)\,dR.
\end{aligned}
\tag{5.11}
\end{equation}
\NSPage{050}%
These are functions of \NSi{NS-F-P050-001}{\eta}, with every lower-order factor held fixed. We write
\NSFormula{NS-F-P050-002}{display}{none}
\[
\mathbf m_n=(m_{n,1},\ldots,m_{n,5}),
\qquad
\mathbf P_n=\left(\phi_n,U_n,\frac{V_n}{X},\Pi_n,\frac{F_n}{X}\right).
\]
Radial supports below are uniform over \NSi{NS-F-P050-003}{|\eta|\leq 1}. For the stress estimates, we use the flat
weight \NSi{NS-F-P050-004}{\zeta} of Theorem~\ref{thm:4.6} (Item~\textup{(iv)}) and set
\NSFormula{NS-F-P050-005}{display}{none}
\[
\delta(X)=\min\{1,\log(X/X_a),\log(X_b/X)\}
\qquad (X_a<X<X_b).
\]

\begin{lemma}\label{lem:5.2}
There are radii \NSi{NS-F-P050-006}{X_a<X_-<X_+<X_b}, independent of \NSi{NS-F-P050-007}{n}, for the following induction.
Fix \NSi{NS-F-P050-008}{n\geq 1}. For each \NSi{NS-F-P050-009}{1\leq j<n}, assume that \NSi{NS-F-P050-010}{\mathbf P_j} is smooth, satisfies
Equations~\eqref{eq:5.2} and~\eqref{eq:5.5} globally, and obeys \eqref{eq:5.12} with
\NSi{NS-F-P050-011}{n=j}. Assume also that these profiles agree with their inner solutions on
\NSi{NS-F-P050-012}{[0,X_-]} and retain the parameter regularity of Lemma~\ref{lem:5.1} throughout
\NSi{NS-F-P050-013}{[0,a^2]}. For \NSi{NS-F-P050-014}{n=1}, there are no lower positive-order hypotheses. Then the inner
order-\NSi{NS-F-P050-015}{n} solution of Lemma~\ref{lem:5.1} has a smooth extension to
\NSi{NS-F-P050-016}{X\geq0}, \NSi{NS-F-P050-017}{|\eta|\leq1}, agreeing with it on
\NSi{NS-F-P050-018}{[0,X_-]} and retaining its parameter regularity throughout
\NSi{NS-F-P050-019}{[0,a^2]}, satisfying Equations~\eqref{eq:5.2} and~\eqref{eq:5.5} globally, and obeying
\NSFormula{NS-F-P050-020}{display}{5.12}
\begin{equation}\label{eq:5.12}
\mathbf m_n(\eta)=0\quad (|\eta|\leq1),
\qquad
\supp_X\mathbf P_n\subset[0,X_+].
\tag{5.12}
\end{equation}
Both \NSi{NS-F-P050-021}{V_n/X} and \NSi{NS-F-P050-022}{F_n/X} extend smoothly to
\NSi{NS-F-P050-023}{X=0}. The stress defined by \eqref{eq:5.9} satisfies
\NSFormula{NS-F-P050-024}{display}{none}
\[
\supp_X\mathcal T_1\subset[X_-,X_b],
\qquad
\supp_X\mathcal T_n\subset[X_-,X_+]\quad(n\geq2),
\]
and, for every fixed profile derivative \NSi{NS-F-P050-025}{\partial^I},
\NSFormula{NS-F-P050-026}{display}{5.13}
\begin{equation}\label{eq:5.13}
|\partial^I\mathcal T_n(X,\eta)|
\leq C_{n,I}\zeta(X)\delta(X)^{-N_{n,I}}
\qquad (X_a<X<X_b,\ |\eta|\leq1).
\tag{5.13}
\end{equation}
Here \NSi{NS-F-P050-027}{C_{n,I},N_{n,I}} may depend on the order and derivative; for
\NSi{NS-F-P050-028}{n\geq2} one may take \NSi{NS-F-P050-029}{N_{n,I}=0}.
\end{lemma}

\begin{proof}
\textbf{Step 1: extend the current profiles while preserving the inner solution.}
We preserve the inner solution on \NSi{NS-F-P050-030}{[0,X_-]}, then choose the extension so that
\NSi{NS-F-P050-031}{\mathbf m_n(\eta)=0} for every \NSi{NS-F-P050-032}{|\eta|\leq1}. The conditions
\NSi{NS-F-P050-033}{m_{n,1}=m_{n,3}=0} remove the exterior streamfunction and pressure terms. The conditions
\NSi{NS-F-P050-034}{m_{n,1}=m_{n,2}=0} cancel the integrated time and viscosity terms;
\NSi{NS-F-P050-035}{m_{n,4}=m_{n,5}=0} cancel the integrated momentum fluxes, including the pressure
contribution. Steps 3 and 4 verify these implications.

Let \NSi{NS-F-P050-036}{(\phi_n^{\mathrm{in}},U_n^{\mathrm{in}},V_n^{\mathrm{in}},\Pi_n^{\mathrm{in}})} be the inner profiles, defined for
\NSi{NS-F-P050-037}{0\leq X\leq a^2} by Lemma~\ref{lem:5.1}. Let \NSi{NS-F-P050-038}{\mathcal I_{\mathrm{pos}}} be the reserved
positive-order patch of Theorem~\ref{thm:4.6} (Item~\textup{(vi)}). We choose \NSi{NS-F-P050-039}{X_-} in the
inner collar of Theorem~\ref{thm:4.6}\textup{(i)}, where the stress cone inequalities have a uniform positive
margin, and fix, independently of \NSi{NS-F-P050-040}{n},
\NSFormula{NS-F-P050-041}{display}{none}
\[
X_-<X_{\mathrm{keep}}<X_{\mathrm{cut}}<a^2<\inf\mathcal I_{\mathrm{pos}},
\qquad
\kappa\in C^\infty([0,\infty)),
\qquad
0\leq\kappa\leq1,
\]
\NSFormula{NS-F-P050-042}{display}{none}
\[
\kappa(X)=1\quad(0\leq X\leq X_{\mathrm{keep}}),
\qquad
\kappa(X)=0\quad(X\geq X_{\mathrm{cut}}).
\]
The initial extensions are
\NSFormula{NS-F-P050-043}{display}{none}
\[
\widetilde\phi_n=\kappa\phi_n^{\mathrm{in}},
\qquad
\widetilde U_n=\kappa U_n^{\mathrm{in}},
\qquad
\widetilde E_n=\frac{R}{\mathsf C}\widetilde\phi_n,
\qquad
R=\sqrt{2X},
\]
with the products extended by zero past \NSi{NS-F-P050-044}{X=a^2}.

In the \NSi{NS-F-P050-045}{R} coordinate, the reserved patch \NSi{NS-F-P050-046}{\mathcal I_{\mathrm{pos}}} is
\NSi{NS-F-P050-047}{\mathcal J_{\mathrm{pos}}=\{R>0:R^2/2\in\mathcal I_{\mathrm{pos}}\}}. We choose nonnegative smooth bumps
\NSFormula{NS-F-P050-048}{display}{none}
\[
b_1^U,b_2^U,b_1^E,b_2^E,b_3^E\in C_c^\infty(\mathcal J_{\mathrm{pos}}),
\qquad
\int_0^\infty b_j^U\,dR=1\quad(j=1,2),
\qquad
\int_0^\infty b_j^E\,dR=1\quad(j=1,2,3),
\]
\NSPage{051}%
with mutually disjoint ordered supports, fixed across orders. For five unknown scalar functions
\NSi{NS-F-P051-001}{\alpha_{n,j}(\eta),\beta_{n,j}(\eta)}, we set
\NSFormula{NS-F-P051-002}{display}{5.14}
\begin{equation}\label{eq:5.14}
\begin{aligned}
U_n(X,\eta)&=\widetilde U_n(X,\eta)+\sum_{j=1}^{2}\alpha_{n,j}(\eta)b_j^U(R),\\
E_n(X,\eta)&=\widetilde E_n(X,\eta)+\sum_{j=1}^{3}\beta_{n,j}(\eta)b_j^E(R),
\qquad
\phi_n=\frac{\mathsf C}{R}E_n.
\end{aligned}
\tag{5.14}
\end{equation}
At \NSi{NS-F-P051-003}{R=0}, the quotient defining \NSi{NS-F-P051-004}{\phi_n} is its smooth extension: near the axis the bumps
vanish and \NSi{NS-F-P051-005}{\phi_n=\widetilde\phi_n}. For each choice of these five coefficient functions, we define the remaining
profiles by
\NSFormula{NS-F-P051-006}{display}{5.15}
\begin{equation}\label{eq:5.15}
\begin{aligned}
F_n(X,\eta)&=\int_0^XU_n(x,\eta)\,dx,\\
V_n(X,\eta)&=-\int_0^X(Z_{-A,n}U_n)(x,\eta)\,dx,\\
\Pi_n(X,\eta)&=\int_0^X\left[\mathsf C^{-2}\sum_{i+j=n}\phi_i\phi_j-\frac{\Omega_{n-1}}{2x}\right](x,\eta)\,dx.
\end{aligned}
\tag{5.15}
\end{equation}
All lower-order profiles and \NSi{NS-F-P051-007}{\Omega_{n-1}} remain fixed in these formulas. The quotient
\NSi{NS-F-P051-008}{\Omega_{n-1}/x} is smooth at the axis, as established after \eqref{eq:5.6}.
Equations~\eqref{eq:5.2} and~\eqref{eq:5.5} hold globally by construction. On
\NSi{NS-F-P051-009}{[0,X_-]}, the cutoff is one and the bumps vanish, so forward integration from the zero axis values gives
\NSFormula{NS-F-P051-010}{display}{none}
\[
(\phi_n,U_n,V_n,\Pi_n)
=(\phi_n^{\mathrm{in}},U_n^{\mathrm{in}},V_n^{\mathrm{in}},\Pi_n^{\mathrm{in}}).
\]
The radial cutoff is independent of \NSi{NS-F-P051-011}{\eta}, and the bumps vanish on
\NSi{NS-F-P051-012}{[0,a^2]}. The reconstructed profiles therefore retain the parameter regularity of
Lemma~\ref{lem:5.1} on that whole interval, as required at the next order.

\textbf{Step 2: solve the five moment equations.} On \NSi{NS-F-P051-013}{\mathcal I_{\mathrm{pos}}},
\NSFormula{NS-F-P051-014}{display}{none}
\[
U_0=0,
\qquad
E_0=e_*f(\eta)R^{-1-2\lambda},
\qquad
e_*>0,
\qquad
\lambda>0,
\]
where \NSi{NS-F-P051-015}{f} is smooth and bounded away from zero. With the lower-order profiles fixed, all five moments are
affine in the five unknown coefficients. For the fifth moment, the pressure equation gives
\NSFormula{NS-F-P051-016}{display}{none}
\[
m_{n,5}=\int_0^\infty\left(
\sum_{i+j=n}U_iU_j
-\frac12\sum_{i+j=n}E_iE_j
+\frac12\Omega_{n-1}
\right)\,dX.
\]
The term \NSi{NS-F-P051-017}{\Omega_{n-1}} depends only on lower orders and is unchanged by the five coefficients. Using
\NSi{NS-F-P051-018}{R^2=2X}, the pressure equation also gives
\NSFormula{NS-F-P051-019}{display}{none}
\[
\partial_R\Pi_n
=\frac{2E_0E_n}{R}
+\frac1R\left(\sum_{i=1}^{n-1}E_iE_{n-i}-\Omega_{n-1}\right).
\]
We define the constant moment matrices
\NSFormula{NS-F-P051-020}{display}{none}
\[
\begin{aligned}
(B_U)_{ij}&=\int_0^\infty R^{p_i}b_j^U(R)\,dR,
& (p_1,p_2)&=(1,1-2\lambda),\\
(B_E)_{ij}&=\int_0^\infty R^{s_i}b_j^E(R)\,dR,
& (s_1,s_2,s_3)&=(2,-2-2\lambda,-2\lambda).
\end{aligned}
\]
\NSPage{052}%
Here \NSi{NS-F-P052-001}{B_U} is \NSi{NS-F-P052-002}{2\times2} and \NSi{NS-F-P052-003}{B_E} is
\NSi{NS-F-P052-004}{3\times3}. Let \NSi{NS-F-P052-005}{m_{n,j}^0} be the moments obtained by setting all five coefficients in
\eqref{eq:5.14} to zero and reconstructing the profiles by \eqref{eq:5.15}. We put
\NSFormula{NS-F-P052-006}{display}{none}
\[
\mathbf d_{U,n}
=\begin{pmatrix}
m_{n,1}^0\\[2pt]
m_{n,4}^0/(e_*f)
\end{pmatrix},
\qquad
\mathbf d_{E,n}
=\begin{pmatrix}
m_{n,2}^0\\[2pt]
m_{n,3}^0/(2e_*f)\\[2pt]
-m_{n,5}^0/(e_*f)
\end{pmatrix}.
\]
Writing \NSi{NS-F-P052-007}{\boldsymbol\alpha_n=(\alpha_{n,1},\alpha_{n,2})^{\mathsf T}} and
\NSi{NS-F-P052-008}{\boldsymbol\beta_n=(\beta_{n,1},\beta_{n,2},\beta_{n,3})^{\mathsf T}}, the five conditions are exactly
\NSFormula{NS-F-P052-009}{display}{none}
\[
\mathbf m_n=0
\quad\Longleftrightarrow\quad
B_U\boldsymbol\alpha_n=-\mathbf d_{U,n},
\qquad
B_E\boldsymbol\beta_n=-\mathbf d_{E,n}.
\]
Lemma~\ref{lem:A.1} applies because the exponents in each block are distinct and the bump supports are ordered.
Both matrices are therefore invertible, and the required coefficients are
\NSFormula{NS-F-P052-010}{display}{5.16}
\begin{equation}\label{eq:5.16}
\boldsymbol\alpha_n=-B_U^{-1}\mathbf d_{U,n},
\qquad
\boldsymbol\beta_n=-B_E^{-1}\mathbf d_{E,n}.
\tag{5.16}
\end{equation}
The discrepancies are smooth functions of \NSi{NS-F-P052-011}{\eta}, and all fixed derivatives of
\NSi{NS-F-P052-012}{1/f} are bounded on \NSi{NS-F-P052-013}{[-1,1]}. Thus these coefficients are smooth for discrepancies of
any finite size. This proves \NSi{NS-F-P052-014}{\mathbf m_n=0}. The support assertion in \eqref{eq:5.12} remains to be checked.

\textbf{Step 3: close the velocity and pressure outside the correction region.}
We choose \NSi{NS-F-P052-015}{X_+} beyond the leading axial-velocity perturbation \NSi{NS-F-P052-016}{U_0} and every
positive-order patch, and before the terminal outer collar. The first moment gives
\NSFormula{NS-F-P052-017}{display}{none}
\[
F_n(X,\eta)=\int_0^\infty U_n(x,\eta)\,dx=m_{n,1}(\eta)=0
\qquad (X\geq X_+).
\]
Hence \NSi{NS-F-P052-018}{\mathcal A_X(U_n)=V_n=0} there at positive order, while the leading axial-flux moment correction
gives \NSi{NS-F-P052-019}{U_0=V_0=0} there as well. Each term of every \NSi{NS-F-P052-020}{\Omega_k} contains a
\NSi{NS-F-P052-021}{V} factor or a derivative of a \NSi{NS-F-P052-022}{V} field, and hence vanishes on this same exterior
region. At positive order each product \NSi{NS-F-P052-023}{\phi_i\phi_j}, \NSi{NS-F-P052-024}{i+j=n}, contains a
positive-order factor. Thus \NSi{NS-F-P052-025}{\partial_X\Pi_n=0} beyond \NSi{NS-F-P052-026}{X_+}, and the third moment gives
\NSFormula{NS-F-P052-027}{display}{none}
\[
\Pi_n(X,\eta)
=\Pi_n(0,\eta)+\int_0^\infty\partial_R\Pi_n(R,\eta)\,dR
=m_{n,3}(\eta)=0
\qquad (X\geq X_+).
\]
Thus \NSi{NS-F-P052-028}{\phi_n,U_n,V_n,\Pi_n,F_n} all vanish for \NSi{NS-F-P052-029}{X\geq X_+}. Their smoothness follows
from the cutoff ansatz, the finite solve \eqref{eq:5.16}, and the integrals \eqref{eq:5.15}. At the axis,
\NSFormula{NS-F-P052-030}{display}{none}
\[
\frac{F_n(X,\eta)}{X}=\int_0^1U_n(tX,\eta)\,dt
\]
is smooth, and \eqref{eq:5.2} gives the same conclusion for \NSi{NS-F-P052-031}{V_n/X}. Compactness of
\NSi{NS-F-P052-032}{[0,X_+]\times[-1,1]} now yields the coefficient bounds
\NSFormula{NS-F-P052-033}{display}{5.17}
\begin{equation}\label{eq:5.17}
\max_{k+\ell\leq m}\ \sup_{X\geq0,\,|\eta|\leq1}
|\partial_X^k\partial_\eta^\ell\mathbf P_n(X,\eta)|
\leq C_{n,m}<\infty
\qquad(n\geq1,\ m\geq0).
\tag{5.17}
\end{equation}
The support radius is common to all orders; these constants may grow with \NSi{NS-F-P052-034}{n} and
\NSi{NS-F-P052-035}{m}.

On the reserved interval \NSi{NS-F-P052-036}{\mathcal I_{\mathrm{mean}}} for later angular-mean corrections, the higher-order
velocity coefficients vanish identically. The support of the inner cutoff and the patch
\NSi{NS-F-P052-037}{\mathcal I_{\mathrm{pos}}} both lie to the left of \NSi{NS-F-P052-038}{\mathcal I_{\mathrm{mean}}}. Thus
\NSi{NS-F-P052-039}{E_n=U_n=0} there, and the first moment gives \NSi{NS-F-P052-040}{F_n=0} already beyond the
positive-order patch; \eqref{eq:5.2} then gives \NSi{NS-F-P052-041}{V_n=0}. In particular,
\NSFormula{NS-F-P052-042}{display}{5.18}
\begin{equation}\label{eq:5.18}
E_n=U_n=F_n=V_n=0
\quad\text{on the reserved mean patch, for every }n\geq1.
\tag{5.18}
\end{equation}
Placing \NSi{NS-F-P052-043}{X_+} beyond the leading axial-velocity perturbation is essential for the pressure support:
\NSi{NS-F-P052-044}{\Omega_0} can remain nonzero beyond the positive-order patches.
\NSPage{053}%
\textbf{Step 4: cancel the total tangential residual integrals.}
The inner equations make both positive-order stress components zero on
\NSi{NS-F-P053-001}{[0,X_-]}. To prove that they also vanish beyond the outer radius, we use the conservative forms of the
tangential residual for any smooth axisymmetric divergence-free field:
\NSFormula{NS-F-P053-002}{display}{none}
\[
\begin{aligned}
r^2\mathscr R_\theta
&=\partial_t(r^2u_\theta)
  +\partial_r(r^2u_ru_\theta)
  +\partial_z(r^2u_zu_\theta)
  -\partial_r(r^2\partial_ru_\theta-ru_\theta)
  -\partial_{zz}(r^2u_\theta),\\
r\mathscr R_z
&=\partial_t(ru_z)
  +\partial_r(ru_ru_z)
  +\partial_z\bigl(r(u_z^2+p)\bigr)
  -\partial_r(r\partial_ru_z)
  -\partial_{zz}(ru_z).
\end{aligned}
\]
Axis parity and compact positive-order supports make the radial boundary terms vanish. At order
\NSi{NS-F-P053-003}{n}, the remaining angular moments before differentiation are
\NSFormula{NS-F-P053-004}{display}{none}
\[
\begin{aligned}
\int_0^\infty r^2u_{\theta,n}\,dr
&=q^{3/2-A+\lambda_n}\int_0^\infty R^2E_n\,dR,\\
\int_0^\infty r^2\sum_{i+j=n}u_{z,i}u_{\theta,j}\,dr
&=q^{3/2-2A+\lambda_n}\int_0^\infty R^2\sum_{i+j=n}U_iE_j\,dR.
\end{aligned}
\]
They vanish because \NSi{NS-F-P053-005}{m_{n,2}=m_{n,4}=0}. The product power is independent of the split
\NSi{NS-F-P053-006}{i+j=n}, so this cancellation holds before taking derivatives in the physical variables
\NSi{NS-F-P053-007}{z,t}. The lower angular moment used by axial viscosity vanishes because
\NSi{NS-F-P053-008}{m_{n-1,2}=0} when \NSi{NS-F-P053-009}{n\geq2}. When \NSi{NS-F-P053-010}{n=1}, the compensated
leading profile instead supplies the convergent identity
\NSFormula{NS-F-P053-011}{display}{5.19}
\begin{equation}\label{eq:5.19}
\int_0^\infty r^2\bigl(u_{\theta,0}(r,z,t)-P(r)\bigr)\,dr=0,
\tag{5.19}
\end{equation}
where \NSi{NS-F-P053-012}{P} is the fixed pure-power exterior of \eqref{eq:A.45}, so
\NSi{NS-F-P053-013}{\partial_z^2P(r)=0}. Differentiation under the integral is justified: beyond the varying terminal collar
\NSi{NS-F-P053-014}{u_{\theta,0}=K(r,t)} is independent of \NSi{NS-F-P053-015}{z}, while the difference moment itself converges.
Thus the order-one angular viscosity term also has zero total moment.

The axial moments are
\NSFormula{NS-F-P053-016}{display}{5.20}
\begin{equation}\label{eq:5.20}
\begin{aligned}
\int_0^\infty ru_{z,n}\,dr
&=q^{1-A+\lambda_n}\int_0^\infty RU_n\,dR,\\
\int_0^\infty r\left(\sum_{i+j=n}u_{z,i}u_{z,j}+p_n\right)\,dr
&=q^{1-2A+\lambda_n}\int_0^\infty R\left(\sum_{i+j=n}U_iU_j+\Pi_n\right)\,dR.
\end{aligned}
\tag{5.20}
\end{equation}
For the pressure contribution, integration by parts gives
\NSFormula{NS-F-P053-017}{display}{none}
\[
\int_0^\infty R\Pi_n\,dR=-\frac12\int_0^\infty R^2\partial_R\Pi_n\,dR.
\]
The identities \NSi{NS-F-P053-018}{m_{n,1}=m_{n,5}=0} therefore make both integrals in \eqref{eq:5.20} vanish.
The axial-viscosity term uses the order-\NSi{NS-F-P053-019}{n-1} axial-flux integral, which also vanishes, including at order
zero. Integrating the conservative equations now gives, order by order,
\NSFormula{NS-F-P053-020}{display}{5.21}
\begin{equation}\label{eq:5.21}
\begin{aligned}
\int_0^\infty r^2\mathscr R_{\theta,n}\,dr
&=\partial_t\int_0^\infty r^2u_{\theta,n}\,dr
 +\partial_z\int_0^\infty r^2\sum_{i+j=n}u_{z,i}u_{\theta,j}\,dr
 -\partial_{zz}\int_0^\infty r^2u_{\theta,n-1}\,dr
 =0,\\
\int_0^\infty r\mathscr R_{z,n}\,dr
&=\partial_t\int_0^\infty ru_{z,n}\,dr
 +\partial_z\int_0^\infty r\left(\sum_{i+j=n}u_{z,i}u_{z,j}+p_n\right)\,dr
 -\partial_{zz}\int_0^\infty ru_{z,n-1}\,dr
 =0.
\end{aligned}
\tag{5.21}
\end{equation}
Here \NSi{NS-F-P053-021}{\mathscr R_{j,n}=q^{-A-1+\lambda_n}\mathfrak r_{j,n}}, for
\NSi{NS-F-P053-022}{j\in\{\theta,z\}}, denotes the order-\NSi{NS-F-P053-023}{n} tangential residual in physical variables. In
\eqref{eq:5.21} at \NSi{NS-F-P053-024}{n=1}, the final integral means the convergent integral of
\NSi{NS-F-P053-025}{u_{\theta,0}-P} from
\NSPage{054}%
\eqref{eq:5.19}; differentiation applies to this convergent subtracted integral. Removing the common physical powers yields
\NSFormula{NS-F-P054-001}{display}{none}
\[
\int_0^\infty R^2\mathfrak r_{\theta,n}\,dR=0,
\qquad
\int_0^\infty R\mathfrak r_{z,n}\,dR=0.
\]
Consequently \eqref{eq:5.9} can also be written as
\NSFormula{NS-F-P054-002}{display}{none}
\[
\begin{aligned}
\mathcal T_{n,\theta}(R)
&=-R^{-2}\int_0^R\varrho^2\mathfrak r_{\theta,n}(\varrho)\,d\varrho
 =R^{-2}\int_R^\infty\varrho^2\mathfrak r_{\theta,n}(\varrho)\,d\varrho,\\
\mathcal T_{n,z}(R)
&=-R^{-1}\int_0^R\varrho\mathfrak r_{z,n}(\varrho)\,d\varrho
 =R^{-1}\int_R^\infty\varrho\mathfrak r_{z,n}(\varrho)\,d\varrho.
\end{aligned}
\]
This representation applies to the signed higher-order stress.

\textbf{Step 5: locate the stress support and bound the order-one outer term.}
For \NSi{NS-F-P054-003}{n\geq2} the residual itself vanishes beyond \NSi{NS-F-P054-004}{X_+}: its same-order terms have a
positive-order factor and its axial viscosity acts on order \NSi{NS-F-P054-005}{n-1>0}. The backward representation proves
\NSi{NS-F-P054-006}{\supp\mathcal T_n\subset[X_-,X_+]}. At \NSi{NS-F-P054-007}{n=1} the only remaining exterior source is
\NSi{NS-F-P054-008}{-\partial_{zz}u_{\theta,0}} in the angular equation. Beyond \NSi{NS-F-P054-009}{X_b}, the leading profile
is the \NSi{NS-F-P054-010}{z}-independent heat field, so \NSi{NS-F-P054-011}{\mathcal T_1} is still supported in the closed active
annulus. For the quantitative edge estimate we restrict further to the terminal collar where the exterior heat-flow cutoff
equals one. There \NSi{NS-F-P054-012}{u_{\theta,0}=K(r,t)f_o}, where
\NSi{NS-F-P054-013}{K(r,t)=c_\infty s^{-A}\mathcal H(2(1-t)/s)} is the physical heat field of \eqref{eq:4.29}.
The smooth terminal multiplier \NSi{NS-F-P054-014}{f_o} is constant on its initial plateau and equals one for
\NSi{NS-F-P054-015}{X\geq X_b}; its exact construction is \eqref{eq:A.12}, with a translation of the logarithmic radial
coordinate. In the local coordinate \NSi{NS-F-P054-016}{y=\log X} it has
\NSi{NS-F-P054-017}{1-f_o=e^{-4/\delta_b^2}} times a smooth factor, with
\NSi{NS-F-P054-018}{\delta_b=\log(X_b/X)}. Direct differentiation gives
\NSFormula{NS-F-P054-019}{display}{none}
\[
a_z(\eta)=-\frac{2\eta}{L},
\qquad
b_z(\eta)=\frac{-2D\eta a_z+d\partial_\eta a_z}{L},
\qquad
\partial_{zz}(Kf_o)=Kq^{-2D}\left(a_z^2\partial_{yy}f_o+b_z\partial_yf_o\right).
\]
All fixed \NSi{NS-F-P054-020}{\eta}-derivatives of \NSi{NS-F-P054-021}{a_z,b_z} are bounded on
\NSi{NS-F-P054-022}{[-1,1]}. The derivatives \NSi{NS-F-P054-023}{\partial_yf_o} and
\NSi{NS-F-P054-024}{\partial_{yy}f_o} are bounded respectively by
\NSi{NS-F-P054-025}{Ce^{-4/\delta_b^2}\delta_b^{-3}} and
\NSi{NS-F-P054-026}{Ce^{-4/\delta_b^2}\delta_b^{-6}}. Backward integration gains three powers of
\NSi{NS-F-P054-027}{\delta_b} by Lemma~\ref{lem:A.9}, applied with
\NSi{NS-F-P054-028}{c=4} and \NSi{NS-F-P054-029}{j=3,6}. The radial weights and the conversion between
\NSi{NS-F-P054-030}{R} and \NSi{NS-F-P054-031}{\delta_b} are bounded smooth factors on this fixed collar. After factoring out
the stress power \NSi{NS-F-P054-032}{q^{-A-1/2+\lambda_1}}, Lemma~\ref{lem:A.9} gives the pointwise estimate
\NSFormula{NS-F-P054-033}{display}{5.22}
\begin{equation}\label{eq:5.22}
|\mathcal T_1|\leq Ce^{-4/\delta_b^2}\delta_b^{-3},
\qquad
|\partial^I\mathcal T_1|\leq C_Ie^{-4/\delta_b^2}\delta_b^{-N_I}.
\tag{5.22}
\end{equation}
The second estimate in \eqref{eq:5.22} follows by differentiating the backward integral; each derivative adds only a finite
inverse power of \NSi{NS-F-P054-034}{\delta_b}. Near this edge the inner factor of \NSi{NS-F-P054-035}{\zeta} is bounded
above and below by positive constants. This proves \eqref{eq:5.13} for \NSi{NS-F-P054-036}{n=1}; the identity
\NSi{NS-F-P054-037}{\mathcal T_1=0} on \NSi{NS-F-P054-038}{[0,X_-]} gives the estimate near the inner edge. For
\NSi{NS-F-P054-039}{n\geq2} the common compact interior support makes every fixed derivative of
\NSi{NS-F-P054-040}{\mathcal T_n/\zeta} bounded.
\end{proof}

\subsection{The coefficient induction and finite residual}\label{sec:5.3}
The preceding lemmas construct and extend one coefficient from the completed lower orders. Lemma~\ref{lem:5.1} solves the
inner equations Equations~\eqref{eq:5.2} to~\eqref{eq:5.6}, and Lemma~\ref{lem:5.2} supplies the extension and support
properties while retaining those equations on \NSi{NS-F-P054-041}{0\leq X\leq X_-}. In this subsection we apply these lemmas
inductively and estimate the residual of each finite truncation in \eqref{eq:5.25}, including its physical derivatives. The gain
in decay must grow with the truncation order, while the exponent loss caused by physical differentiation stays independent
of that order, as quantified below in \eqref{eq:5.26}.
